\section{Introduction}
\label{sec:introduction}

Automated recognition of facial race and gender attributes plays a crucial role in digital identity governance, biometric access control, digital forensics, and intelligent interface personalization~\cite{ref1}, \cite{ref2}. Numerous independent studies and evaluations by the National Institute of Standards and Technology (NIST) demonstrate that commercial computer vision models exhibit performance disparities that disadvantage specific minority groups, particularly dark-skinned women~\cite{ref3}, \cite{ref4}. This accuracy gap stems from imbalanced training data distribution characteristics as well as visual feature representations that lack robustness against phenotypic variations and environmental illumination~\cite{ref5}. The limitations of conventional systems in handling such visual variations highlight the need for novel representation approaches that model race and gender interactions in an integrated manner.

A common practice in computer vision literature separates race and gender classification into two independent, standalone computational models~\cite{ref6}. Such an isolated approach fails to capture interaction dynamics and errors at the intersection of specific demographic subgroups, such as Black women or Asian women~\cite{ref7}, \cite{ref8}. Consequently, classification performance degradation on these specific demographic intersections is often overlooked by conventional aggregate accuracy metric audits~\cite{ref9}. An integrated intersectional formulation mapping six compound classes (3 Race $\times$ 2 Gender) under balanced evaluation conditions enables a fair and transparent performance audit. This intersectional classification modeling demands rich visual representations capable of separating decision boundaries between phenotypically overlapping subgroups~\cite{ref10}.

Facial attribute recognition methodologies have evolved from classical handcrafted features toward Convolutional Neural Network (CNN) and Vision Transformer (ViT) architectures. In the CNN era, prior research explored demographic feature extraction utilizing middle-face regions~\cite{ref11}, as well as multi-scale multi-task frameworks for age and gender estimation~\cite{ref12}. In addition, deep CNN models were developed to enhance the accuracy of facial image-based ethnicity recognition~\cite{ref13}. The ViT architecture subsequently emerged, offering the advantage of modeling global spatial relationships across image patches through the Multi-Head Self-Attention (MHSA) mechanism~\cite{ref14}. Synthetic generative representation approaches have also been explored to mitigate demographic disparities across gender and racial groups~\cite{ref15}. Recent developments include the integration of parallel attention modules~\cite{ref16}, as well as the utilization of ViTs for facial gender classification~\cite{ref17}. Furthermore, exploration of multi-axis transformer architectures has been proposed to enrich the diversity of hierarchical visual representation modeling~\cite{ref18}. Combining cross-domain ViT representations has begun to be pioneered through the integration of dual face and emotion features~\cite{ref19}, as well as multi-domain representation fusion for fair demographic attribute recognition~\cite{ref20}. Despite these rapid advancements, the literature review identifies several methodological challenges that still require integrated empirical investigation.

The limitation of single-domain representations constitutes a primary methodological gap, as static facial biometrics are vulnerable to distortions caused by dynamic affective expressions~\cite{ref11}, \cite{ref17}. Moreover, facial recognition systems are sensitive to biological aging variations when age morphological features are neglected~\cite{ref12}, \cite{ref18}. Another critical gap is the threat of identity leakage, wherein conventional image-level random partitioning splits photographs of the same subject across training and testing sets, thereby inducing identity memorization. Furthermore, prior studies generally lack formal inferential significance verification~\cite{ref13}, \cite{ref19}. Practical computational efficiency profiling of downstream classifiers has also not been comprehensively explored~\cite{ref20}. To bridge these gaps, this study proposes a tri-domain ViT fusion under a rigorous evaluation protocol.

To address these gaps, this study proposes a tri-domain ViT fusion that combines static biometric geometry (ViT-Face), dynamic affective expressions (ViT-Emotion), and biological aging morphology (ViT-Age) through direct concatenation (Face $\oplus$ Emotion $\oplus$ Age: 2,304 dimensions)~\cite{ref21}, \cite{ref22}. Visual latent representations are extracted offline from frozen pre-trained ViT backbones to maintain computational efficiency without additional fine-tuning~\cite{ref23}. This study employs a dual-benchmark framework establishing a strict subject-disjoint protocol (600 unique subjects, 480 training and 120 testing subjects isolated, 0\% leakage) as the primary standard for generalization, while retaining the subject-overlapping Stratified Split evaluation protocol as an ablation study to quantify identity memorization. Comprehensive hyperparameter optimization is conducted across four classical classifier families, namely Random Forest (RF), Gaussian Naive Bayes (GNB), Logistic Regression (LR), and Support Vector Machine (SVM), via 5-Fold StratifiedGroupKFold cross-validation. Modular scaling and Principal Component Analysis (PCA) pipelines are fitted strictly on the training folds in isolation to prevent information leakage~\cite{ref24}.

The primary scientific contributions of this study are summarized in four essential aspects as follows:
\begin{enumerate}
\item Formulation of a seven-feature-scheme ablation framework integrating frozen pre-trained ViT backbones with leakage-free fold-level preprocessing and hyperparameter optimization across four classifier families for intersectional demographic classification.
\item Establishment of an identity-leakage-free evaluation benchmark on the DemogPairs dataset and a comparative data leakage ablation study to quantify the effects of facial identity memorization.
\item Paired inferential statistical significance evaluation validating the significant superiority of multi-domain representations over single-domain baselines, demonstrating predictive parity between the top-performing linear and non-linear models, and uncovering linear age feature complementarity versus non-linear asymptotic representation saturation.
\item Presentation of a comprehensive computational efficiency profile encompassing ViT extraction time, training cost, inference latency, and memory footprint, realizing ultra-lightweight downstream classification heads for practical deployment.
\end{enumerate}

The remainder of this article is organized as follows: Section~\ref{sec:related_works} reviews the literature synthesis on facial demographic recognition and transformer architectures. Section~\ref{sec:materials_and_methods} details the dataset specifications, multi-domain ViT feature extraction methodology, and classifier optimization pipeline. Section~\ref{sec:results_and_discussion} presents experimental evaluation results, data leakage ablation studies, inferential testing, computational efficiency profiles, and intersectional fairness analysis. Finally, Section~\ref{sec:conclusion} summarizes key conclusions, methodological limitations, and recommended directions for future research.
